\documentclass[twocolumn,superscriptaddress,unsortedaddress,nofootinbib,amsmath,amssymb,
aps,floatfix]{revtex4-2}
\usepackage{graphicx,tikz} 
\usetikzlibrary{arrows.meta,decorations.pathreplacing,calc,positioning}
\usepackage[mode=buildmissing]{standalone} 
\usepackage{dcolumn}
\usepackage{bm}
\usepackage{physics}
\usepackage{xcolor}
\usepackage{siunitx}
\usepackage{soul}
\usepackage{comment}
\usepackage[bb=boondox]{mathalfa}
\usetikzlibrary{external}
\tikzexternalize[prefix=figures/]
\DeclareMathOperator{\sinc}{sinc}
\newcommand{\may}[1]{\textcolor{orange}{#1}}
\begin{document}
\title{Unitary Coupled Cluster theory on Quantum Computers with Parametric Matrix Models}
\author{Patrick Cook and Danny Jammooa}
\affiliation{Facility for Rare Isotope Beams and Department of Physics and Astronomy, Michigan State University, East Lansing, Michigan 48824, USA}
\author{Morten Hjorth-Jensen}
\affiliation{Department of Physics and Center for Computing in Science Education, University of Oslo, N-0316 Oslo, Norway}
\author{Dean Lee}
\affiliation{Facility for Rare Isotope Beams and Department of Physics and Astronomy, Michigan State University, East Lansing, Michigan 48824, USA}
\begin{abstract}
Text to come here
\end{abstract}

%\pacs{02.70.Ss, 31.15.A-, 31.15.bw, 71.15.-m, 73.21.La}


\date{\today}

\maketitle

\section{Introduction} 
Structure of article
\begin{enumerate}
    \item Introduction with motivation
    \item Discussion of the Lipkin model
    \item Setting up time evolution
    \item PMMs and time evolution
    \item Trotterization and CC theory
\end{enumerate}

\section{Mathematical Framework and Entanglement Measures}

\subsection{Hamiltonian and Symmetry.}
The Lipkin--Meshkov--Glick (LMG) model describes $N$ spin-$\tfrac{1}{2}$ particles interacting through collective spin operators.  
For four fermions, the case we consider in the examples  here, each levels has
degeneration $d=4$, leading to different total spin values.  The two
levels have quantum numbers $\sigma=\pm 1$, with the upper level
having $2\sigma=+1$ and energy $\varepsilon_{1}= \varepsilon/2$. The
lower level has $2\sigma=-1$ and energy
$\varepsilon_{2}=-\varepsilon/2$. That is, the lowest single-particle
level has negative spin projection (or spin down), while the upper
level has spin up.  In addition, the substates of each level are
characterized by the quantum numbers $p=1,2,3,4$.



The Hamiltonian of the system is given by
\[
\begin{array}{ll}
\hat{H}=&\hat{H}_{0}+\hat{H}_{1}+\hat{H}_{2}\\
&\\
\hat{H}_{0}=&\frac{1}{2}\varepsilon\sum_{\sigma ,p}\sigma
a_{\sigma,p}^{\dagger}a_{\sigma ,p}\\
&\\
\hat{H}_{1}=&\frac{1}{2}V\sum_{\sigma ,p,p'}
a_{\sigma,p}^{\dagger}a_{\sigma ,p'}^{\dagger}
a_{-\sigma ,p'}a_{-\sigma ,p}\\
&\\
\hat{H}_{2}=&\frac{1}{2}W\sum_{\sigma ,p,p'}
a_{\sigma,p}^{\dagger}a_{-\sigma ,p'}^{\dagger}
a_{\sigma ,p'}a_{-\sigma ,p}\\
&\\
\end{array}
\]
where $V$ and $W$ are constants. The operator 
$H_{1}$ can move pairs of fermions
while $H_{2}$ is a spin-exchange term. The latter
moves a pair of fermions from a state $(p\sigma ,p' -\sigma)$ to a state
$(p-\sigma ,p'\sigma)$.

We are going to rewrite the above Hamiltonian in terms of so-called  quasispin operators
\[
\begin{array}{ll}
\hat{J}_{+}=&\sum_{p}
a_{p+}^{\dagger}a_{p-}\\
&\\
\hat{J}_{-}=&\sum_{p}
a_{p-}^{\dagger}a_{p+}\\
&\\
\hat{J}_{z}=&\frac{1}{2}\sum_{p\sigma}\sigma
a_{p\sigma}^{\dagger}a_{p\sigma}\\
&\\
\hat{J}^{2}=&J_{+}J_{-}+J_{z}^{2}-J_{z}\\
&\\
\end{array}
\]
and the number operator 
\[
\hat{N}=\sum_{p\sigma}
a_{p\sigma}^{\dagger}a_{p\sigma}.
\]

We will now rewrite the Hamiltonian in terms of the above quasi-spin operators and the number operator 
\begin{equation}
N = \sum_{p,\sigma} a_{p\sigma}^\dagger a_{p\sigma}.
label{eq:N}
\end{equation}
Going through each term of the Hamiltonian and using the expressions for the quasi-spin operators we obtain
\begin{equation}
H_0 = \varepsilon J_z, 
\end{equation}
\begin{equation}
H_1 = \frac{1}{2} V \left( J_+^2 + J_-^2 \right),
\end{equation}
and 
\begin{equation}
H_2 = \frac{1}{2} W \left( -N + J_+ J_- + J_- J_+ \right).
\end{equation}

The above expressions can in turn be used to show that the Hamiltonian
commutes with the various quasi-spin operators. This leads to quantum
numbers which are conserved.  Let us first show that $[H,J^2]=0$,
which means that $J$ is a so-called *good* quantum number and that the
total spin is a conserved quantum number.

The matrix elements are given by $\langle J,J_z \vert H \vert J',J_z' \rangle$.
The
Hamiltonian is hermitian and we obtain after all this labor of ours

\begin{equation}
H_{J = 2} =
\begin{bmatrix}
-2\varepsilon & 0 & \sqrt{6}V & 0 & 0 \\
0 & -\varepsilon + 3W & 0 & 3V & 0 \\
\sqrt{6}V & 0 & 4W & 0 & \sqrt{6}V \\
0 & 3V & 0 & \varepsilon + 3W & 0 \\
0 & 0 & \sqrt{6}V & 0 & 2\varepsilon
\end{bmatrix}
label{eq:HJ=2}
\end{equation}




\section{Parametric Matrix Models}

\subsection{Scientific Background and Motivation}

The ability to model, predict, and control the behavior of complex
quantum systems is central to the ongoing second quantum
revolution. Across quantum chemistry, condensed matter, and quantum
technology, the challenge lies in bridging the gap between microscopic
Hamiltonian descriptions and emergent macroscopic
phenomena. Traditional methods—such as exact diagonalization or
coupled-cluster theory—are computationally demanding, often scaling
exponentially with system size. Machine learning (ML) offers a
complementary route by enabling data-driven discovery of patterns and
effective models. However, standard ML approaches typically lack
interpretability, stability, and physical consistency.


A Parametric Matrix Model describes a system (or dataset) through a
matrix whose **entries depend on a set of tunable
parameters. Formally, one writes
\[
M(\theta) = \sum_i \theta_i A_i,
\]
where $A_i$ are fixed basis matrices and $\theta_i$ are parameters to
be learned or optimized.  Such models can represent Hamiltonians,
covariance matrices, kernels, or linear transformations whose
structure reflects known physics, symmetries, or constraints.

Unlike dense neural networks that learn arbitrary mappings, PMMs
impose a low-dimensional, interpretable structure. This makes them
ideal for
\begin{itemize}
\item learning models with few but physically meaningful parameters,
\item ensuring stability and symmetry preservation (e.g., Hermiticity, orthogonality),
\item reducing overfitting and improving **generalization**.
\end{itemize}


Because $M(\theta)$ is parametric and smooth in its parameters, it can
be integrated seamlessly with automatic differentiation frameworks
(e.g., PyTorch, TensorFlow). This allows one to train PMMs end-to-end
as part of a neural architecture—for instance, as a differentiable
layer encoding physical interactions or correlations.

In quantum mechanics, the spectrum and dynamics of a system depend on
its Hamiltonian matrix. PMMs allow ML models to learn maps from
parameters to observables (e.g., energies, eigenvectors, correlation
functions). This has direct implications for:
\begin{itemize}
\item learning effective Hamiltonians,
\item building surrogate models for quantum simulations,
\item predicting **phase transitions** or control parameters.
\end{itemize}

PMMs enable hybrid learning, where parts of the model are learned from
data and parts are constrained by physical principles. For example
\begin{itemize}
\item embedding PMMs in a neural network can encode **unitarity, time-reversal symmetry, or particle-number conservation,
\item combining PMMs with variational learning yields **interpretable latent spaces** corresponding to physically relevant degrees of freedom.
\end{itemize}

The recent work \cite{pmm2025} shows that PMMs can learn physically
consistent effective theories directly from data — in particular, how
parameters influence spectral and dynamical properties.  This bridges
machine learning and quantum many-body physics, allowing one to
\begin{itemize}
\item  infer Hamiltonian landscapes,
\item emulate complex correlated systems, and
\item  guide quantum control or materials design.
\end{itemize}

Parametric Matrix Models are valuable in ML because they
\begin{enumerate}
\item Encode structure, symmetries, and interpretability.
\item Reduce model complexity and improve generalization.
\item Provide differentiable, physics-informed representations.
\item Connect data-driven learning with analytical and physical insight.
\end{enumerate}

In short — they offer a bridge between pure data fitting and
mechanistic understanding, making them ideal for scientific machine
learning and quantum technology applications.



\section{Coupled Cluster Theory (Review)}
In classical many-body quantum chemistry, the \emph{Coupled Cluster} (CC) method uses an exponential ansatz for the wavefunction.  One writes 
\begin{equation}
|\Psi_{\rm CC}\rangle = e^{\hat T}\,|\Phi_0\rangle,
\end{equation}
where $|\Phi_0\rangle$ is a reference (usually Hartree--Fock) determinant, and the cluster operator $\hat T$ is a sum of excitation operators 
\begin{equation}
\hat T = \hat T_1 + \hat T_2 + \cdots + \hat T_N.
\end{equation}

Here $\hat T_1$ generates all single-particle excitations from
occupied to virtual orbitals, $\hat T_2$ generates all double
excitations, and so on.  For example, in second-quantized notation
with creation ($a^\dagger_a$) and annihilation ($a_i$) operators, one
may write
\[
\hat T_1 = \sum_{i,a} t_i^a\,a^\dagger_a a_i,\qquad
\hat T_2 = \frac{1}{4}\sum_{i<j,a<b} t_{ij}^{ab}\,a^\dagger_a a^\dagger_b a_j a_i,
\]
and in general $T_n$ promotes $n$ electrons.  Truncation of the sum yields common CC models: e.g.\ the singles-and-doubles approximation (CCSD) retains only $T_1+T_2$

The exponential form ensures \emph{size-consistency}
(size-extensivity) of the method: non-interacting subsystems have
additive energy because $e^{\hat T}$ factorizes By contrast, a
truncated configuration interaction (CI) expansion lacks this
property.  The nonlinearity of $e^{\hat T}$ also effectively includes
portions of higher excitations (e.g.\ $\hat T_2^2$ yields quadruple
excitations) even when truncating at doubles
which accelerates convergence toward the exact (FCI) solution.  In
practice, the CC energy is obtained by inserting the ansatz into the
Schrödinger equation $H|\Psi_{\rm CC}\rangle = E|\Psi_{\rm CC}\rangle$
and projecting against excited determinants.  Equivalently, one works
with the non-Hermitian similarity-transformed Hamiltonian $\bar H=e^{-\hat T}He^{\hat T}$, solving
\[
\langle\Phi_0|\bar H|\Phi_0\rangle = E_{\rm CC},\quad 
\langle\Phi_{\mu}|\bar H|\Phi_0\rangle = 0,
\]
for all excited determinants $\ket{\Phi_\mu}$ up to the truncation
level.  Because $e^{\hat T}$ is not unitary, the CC energy is not
obtained via a variational minimization (indeed conventional CC is not
variational In fact, CCSD scales as $O(N^6)$ with system size (more
generally, CC with up to $k$-fold excitations scales like $N^{k+2}$).
Despite these costs, CCSD and its perturbatively-augmented variant
CCSD(T) are among the most accurate methods for weakly-correlated
molecular systems.

\subsection{The Unitary Coupled Cluster Ansatz}
The \emph{Unitary Coupled Cluster} (UCC) approach modifies the CC ansatz to make it unitary, suitable for quantum hardware.  In UCC one defines an anti-Hermitian cluster operator $\hat A = \hat T - \hat T^\dagger$ and writes 
\begin{equation}
|\Psi_{\rm UCC}\rangle = e^{\hat T - \hat T^\dagger}\,|\Phi_0\rangle,
\end{equation}
so that $e^{\hat A}$ is manifestly unitary.  For UCCSD, $\hat T$ includes the usual singles and doubles excitation terms, but each excitation amplitude $t_i^a$ is paired with its de-excitation to form $(t_i^a - t_i^{a*})$ in the exponent. 
This ansatz preserves the size-extensivity of CC but ensures unitarity, making it a natural variational form for quantum circuits.

In principle the UCC energy
\[
E=\langle\Phi_0|e^{-(\hat T-\hat
  T^\dagger)} H e^{\hat T-\hat T^\dagger}|\Phi_0\rangle
\]
could be
variationally minimized.  However, classically evaluating this
expectation is intractable for large systems due to the
non-terminating Baker–Campbell–Hausdorff expansion.  Thus UCC is
rarely used in classical CC codes.  On a quantum computer, the unitary
exponentials can be directly implemented, sidestepping the need for
the non-Hermitian similarity transform

Key differences from classical CC include: UCC uses a unitary operator
and hence can be optimized by energy minimization, but the operators
$\hat \tau_i - \hat \tau_i^\dagger$ do not commute in general.  The
resulting UCC ansatz (e.g.\ UCCSD) thus defines a manifold of
wavefunctions different from conventional CCSD.  In practice one often
constructs UCC ansätze by decomposing $e^{\hat T-\hat T^\dagger}$ into
implementable quantum gates (via Trotterization; see below).  In
applications, UCCSD has been extensively studied in quantum chemistry
(as a variational wavefunction and as a benchmark)

Its advantages on quantum devices include the exact preservation of
norm and the ability to define variational parameters directly as gate
angles, at the cost of deeper circuits.

\subsection{Trotterization of the UCC Operator}

The unitary operator $U_{\rm UCC}=e^{\hat T-\hat T^\dagger}$ generally
cannot be implemented in one piece on a quantum circuit, since $\hat
T-\hat T^\dagger = \sum_i (\hat \tau_i - \hat \tau_i^\dagger)$ is a
sum of non-commuting terms (each $\hat \tau_i$ being a fermionic
excitation generator).  Instead, one applies the \emph{Trotter-Suzuki}
(product-formula) decomposition to approximate the exponential of a
sum by a product of exponentials of individual terms.  Formally, for a
Trotter step size $1/t$, one writes

\[
e^{\sum_i A_i} = \lim_{t\to\infty}\Bigl(\prod_i e^{A_i/t}\Bigr)^t.
\]
In first-order (Lie) Trotter approximation ($t=1$), one simply takes the product once:
\begin{equation}\label{eq:trotter1}
e^{\hat T-\hat T^\dagger} \approx \prod_i \exp\!\bigl(\theta_i(\hat\tau_i - \hat\tau_i^\dagger)\bigr),
\end{equation}
where $\theta_i$ are the corresponding cluster amplitudes.
More generally, with $t$ steps one has 
\[
e^{\hat T-\hat T^\dagger} \approx \Bigl[\prod_i e^{(\theta_i/t)(\hat\tau_i - \hat\tau_i^\dagger)}\Bigr]^t + O(1/t),
\]
with an $O(1/t)$ error per step (i.e.\ $O(\Delta^2)$ error in total with step size $\Delta=1/t$).

Physically, this means repeating a sequence of gates $t$ times can
systematically improve the accuracy, at the cost of circuit depth.  In
practice, one often takes $t=1$ or a small integer in order to limit
gate count.

Higher-order Suzuki formulas can reduce the error scaling.  For
example, the second-order symmetric Trotter decomposition splits each
exponential as
\begin{equation}
e^{A+B} \approx e^{A/2} e^{B} e^{A/2},
\end{equation}
with error $O(\Delta^3)$.

Applied to UCCSD, one could interleave half-rotations of one
excitation group, a full rotation of another, then the other
half-rotation back.  Generally, a $k$th-order Suzuki formula reduces
the Trotter error but requires more gate operations (for example, a
second-order sequence of $m$ exponentials yields roughly $2m-1$
exponentials per layer).  In the quantum chemistry VQE context,
first-order (even single-step) Trotterization is commonly used since
variational optimization can partly absorb the Trotter error,
(By contrast, in quantum phase estimation one needs high Trotter order to faithfully simulate real time-evolution.) 

\subsection{UCC Ansatz in the VQE Framework}

The Unitary Coupled Cluster ansatz is widely employed as a trial state
in the Variational Quantum Eigensolver (VQE) for molecular ground
states.

In VQE, one constructs a parameterized unitary circuit
$U(\vec\theta)$, prepares the state $|\Psi(\vec\theta)\rangle =
U(\vec\theta)|\Phi_0\rangle$, and measures its energy
$E(\vec\theta)=\langle\Psi(\vec\theta)|H|\Psi(\vec\theta)\rangle$.  A
classical optimizer then adjusts the parameters $\vec\theta$ to
minimize $E(\vec\theta)$.

Formally, one solves 
\[
E_{\min} = \min_{\vec\theta}\langle\Phi_0|U^\dagger(\vec\theta)\,H\,U(\vec\theta)|\Phi_0\rangle,
\]
and the optimal state is $|\Psi_{\rm opt}\rangle = U(\vec\theta^*)|\Phi_0\rangle$.

In the case of a UCCSD ansatz,
\[
U(\vec\theta)=e^{\hat T(\vec\theta)-\hat T^\dagger(\vec\theta)}$ ensures $|\Psi_{\rm opt}\rangle,
\]
is a unitary-rotated HF state.  Crucially, $U_{\rm UCCSD}$ must
include gates that prepare the reference HF state if it does not
coincide with the all-zero qubit state.

Implementation steps include: first, perform a classical Hartree--Fock
calculation to obtain molecular orbitals and the one- and two-electron
integrals.  Then build the second-quantized Hamiltonian
$H=\sum_{pq}h_{pq}\,a_p^\dagger a_q +
\frac12\sum_{pqrs}h_{pqrs}\,a_p^\dagger a_q^\dagger a_r a_s$ in that
basis.  Also define the excitation operators $\hat\tau_i$
corresponding to chosen single and double excitations.  Next, map the
fermionic operators to qubit Pauli operators (e.g.\ via Jordan–Wigner
or Bravyi–Kitaev transforms)

For example, $a^\dagger_p a_q - a^\dagger_q a_p$ maps to a sum of
Pauli strings.  The mapped UCCSD unitary is then decomposed into
native gates using the Trotter formulas above.  Each factor
$\exp[\theta_i(\hat\tau_i-\hat\tau_i^\dagger)]$ can be further
decomposed into a sequence of single- and two-qubit rotations and
CNOTs using standard techniques


During the VQE loop, the parameter set $\vec\theta$ (initially set
e.g.\ to MP2 amplitudes or zeros) is optimized.  Each iteration
involves: (1) preparing the quantum circuit $U(\vec\theta)$ on the
qubits starting from $|\Phi_0\rangle$, (2) measuring expectation
values of the Hamiltonian terms (requiring repeated state
preparations), and (3) updating $\vec\theta$ on a classical computer
according to a chosen optimization algorithm.  The process repeats
until convergence of the energy.  Because the ansatz is unitary, the
Rayleigh–Ritz variational principle guarantees the energy estimate is
an upper bound to the true ground state.  In practice, optimization
may be performed using gradient-based or gradient-free methods, often
requiring many quantum circuit evaluations.  The quality of the result
depends on the expressiveness of the ansatz and the optimizer’s
ability to navigate potential local minima


\subsection{Algorithmic Implementation}
The implementation of UCCSD within VQE can be summarized as follows 

\begin{algorithm}[H]
\caption{VQE with UCCSD Ansatz}
\begin{algorithmic}[1]
\State \textbf{Classical precomputation:} Perform Hartree--Fock to obtain integrals and reference state $|\Phi_0\rangle$.
\State Build the second-quantized Hamiltonian $\hat H$ and identify excitation operators $\{\hat\tau_i\}$ (singles, doubles, \dots).
\State Choose a mapping (e.g.\ Jordan–Wigner) to convert $\{\hat H,\hat\tau_i\}$ into qubit Pauli operators
\State \textbf{Ansatz construction:} Define parameterized unitary $U(\vec\theta)=\prod_i e^{\theta_i(\hat\tau_i - \hat\tau_i^\dagger)}$ (Trotterized UCCSD).
\State Initialize parameters $\vec\theta$ (e.g.\ from MP2 amplitudes or zeros).
\Repeat
    \State Prepare state $|\Psi(\vec\theta)\rangle = U(\vec\theta)\ket{\Phi_0}$ on the quantum processor.
    \State Measure the energy $E(\vec\theta) = \bra{\Psi(\vec\theta)}\hat H\ket{\Psi(\vec\theta)}$ by sampling (estimating each Pauli term)
    \State Use a classical optimizer to update $\vec\theta \leftarrow \vec\theta - \eta \nabla E(\vec\theta)$ (or other rule) to lower the energy.
\Until{convergence of $E$ or $\vec\theta$ changes small}
\State \textbf{Return:} Optimized parameters $\vec\theta^*$ and energy $E(\vec\theta^*)$ (approximate ground state energy).
\end{algorithmic}
\end{algorithm}

This high-level algorithm illustrates the hybrid nature of VQE.  In practice, efficient gradient evaluation or reduced-measurement schemes may be employed.  The accuracy hinges on including sufficient excitation operators in $U(\vec\theta)$; for UCCSD one includes all single and double excitations with distinct spin/orbital indices.  Diagrammatic or tensor-network derivations are usually not needed on the quantum side since the circuit implicitly handles the entangling structure.  The mapping step (line 3) is often implemented by software tools (e.g.\ OpenFermion, Qiskit) which output a list of Pauli strings.  The exponential of each string can then be translated to gates via standard formulas.

The Unitary Coupled Cluster (UCC) method offers a powerful,
chemically-inspired ansatz for quantum algorithms.  It retains the
accuracy and size-extensivity of classical CC (via the exponential
parametrization) while making the state preparation unitary for
quantum hardware.  By employing Trotter–Suzuki decompositions, the UCC
operator is factorized into implementable gate sequences.  Within VQE,
UCCSD provides a systematic, variational approach to molecular
eigenstates.

Future improvements include adaptive ordering of terms to reduce
circuit depth, higher-order Trotter schemes to lower approximation
error, and use of symmetries to prune the ansatz.  Overall, UCC
exemplifies the synergy of many-body theory with quantum computing,
bridging traditional electronic structure methods and near-term
quantum algorithms.


\section{Results and discussions}


\section{Conclusion}


\end{document}
